\scanpage{662}
\mspage{\textcircled{50}}
% archive sheet number: 661
ça ce sont les exemples qui manquent. La \unsure{clef} d’une compréhension semble bien être dans les questions de factorisation entre ensembles $\Phi\longleftrightarrow\Psi$\unsure{,|.} Il faudrait en faire une étude systématique, mais je recule devant la perspective d’une digression qui me prendrait peut-être encore des semaines !

\noindent\hrulefill\par\noindent\hrulefill

4. \inserted{\uline{Paratopos et}} \uline{L’opération} $M\boxtimes N$.

4.1. \inserted{\uline{Pseudo-topos.} \uline{\unsure{Pts. de première|Pts. du premier} Catégorie}} J’appelle \uline{pseudo-$\unsure{\mathcal U|\mathrm U}$-topos}, ou simplement \uline{pseudo-topos}, une $\mathcal U$-catégorie $M$ stable par $\mathcal U$-limites \deleted{\unsure{projectives}} \inserted{inductives}, et qui admet une \uline{petite} sous-catégorie $M_0$ \deleted{stable par} \uline{strictement génératrice} (i.e. génératrice par épimorphismes stricts).

Alors $M$ admet aussi des $\mathcal U$-limites projectives \inserted{(SGA 4 I 8.12.7)}. [\unsure{Mais} l’exemple de \unsure{(Gr)|(Gr.)}, cf loc cit 8.12.9, montre que $M^{\mathrm o}$ n’est pas forcément, pour \unsure{autant}, un pseudo-$\mathcal U$-topos, faute à $M$ d’admettre une petite \deleted{\unsure{famille}} sous-catégorie strictement \uline{cogénératrice}. Mais le th. 2.6. donne des conditions sur $M$ (\unsure{surtout} du type \unsure{exactitude})

